\documentclass[11pt]{article}
\usepackage{amsmath,amssymb,amsthm}
\usepackage{hyperref}
\usepackage{graphicx}
\usepackage{xcolor}
\usepackage{booktabs}
\usepackage{longtable}
\usepackage{array}
\usepackage{calc}
\usepackage{soul}
\usepackage{geometry}
\providecommand{\tightlist}{\setlength{\itemsep}{0pt}\setlength{\parskip}{0pt}}
\geometry{margin=1in}

% Theorem environments
\newtheorem{theorem}{Theorem}[section]
\newtheorem{lemma}[theorem]{Lemma}
\newtheorem{proposition}[theorem]{Proposition}
\newtheorem{corollary}[theorem]{Corollary}
\newtheorem{conjecture}[theorem]{Conjecture}
\theoremstyle{definition}
\newtheorem{definition}[theorem]{Definition}
\newtheorem{assumption}[theorem]{Postulate}
\theoremstyle{remark}
\newtheorem{remark}[theorem]{Remark}

% Cause-plex notation shortcuts
\newcommand{\CP}{\mathcal{C}}
\newcommand{\Cstar}{\mathcal{C}^*}
\newcommand{\loopspace}{\Omega(\mathcal{C}^*, e_*)}

\title{Cause-Plex and Spacetime: Deriving the Lorentzian Metric from Causal Structure}
\author{Ian Derrington}
\date{2026-03-25}

\begin{document}
\maketitle

\begin{abstract}
Formal derivation of the Lorentzian metric, Lorentz invariance, and energy conservation from cause-plex structure alone. The primitive is the causal event — a state transition with no assumed physics. Spacetime and energy emerge from structural properties of the cause-plex hypergraph.
\end{abstract}

\begin{quote}
This document derives spacetime structure from cause-plex primitives. It
draws on causal set theory (Bombelli, Lee, Meyer, Sorkin 1987),
Malament's theorem (1977), and Wolfram's causal invariance argument.
Where results are established, proofs are given. Where results are at
the research frontier (GR from causal dynamics), status is stated
explicitly.
\end{quote}

\medskip

\section{The Primitive}\label{the-primitive}

\begin{definition}[Causal event]
\label{definition:1-1}
A causal event (\emph{state couplet}) $e$ is an ordered pair $(\mathcal{S}_{\text{in}}, \mathcal{S}_{\text{out}})$ where $\mathcal{S}_{\text{out}}$ is determined by $\mathcal{S}_{\text{in}}$. No physical content is assumed — no energy, no units, no conservation laws.

\end{definition}

\begin{definition}[Cause-plex]
\label{definition:1-2}
The cause-plex $\mathcal{C} = (E, \prec)$ is a set of causal events $E$ equipped with a binary relation $\prec$ ("precedes") satisfying:
\end{definition}

\begin{enumerate}
\def\labelenumi{\arabic{enumi}.}
\tightlist
\item
  \textbf{Irreflexivity:} \(e \not\prec e\) for all \(e \in E\)
\item
  \textbf{Transitivity:} \(e_1 \prec e_2\) and \(e_2 \prec e_3\) implies
  \(e_1 \prec e_3\)
\item
  \textbf{Local finiteness:} for all \(e_1, e_2 \in E\), the set
  \(\{e \in E : e_1 \prec e \prec e_2\}\) is finite
\end{enumerate}

\((E, \prec)\) is a \textbf{locally finite partially ordered set}
(locally finite poset). This is identical to the structure of a causal
set as defined in causal set theory.

\begin{definition}[Spacelike separation]
\label{definition:1-3}
Events $e_1, e_2 \in E$ are *spacelike separated*, written $e_1 \perp e_2$, if neither $e_1 \prec e_2$ nor $e_2 \prec e_1$.

\end{definition}

\begin{definition}[State domain]
\label{definition:1-4}
Each event $e$ has a *state domain* $D(e) \subseteq \mathcal{S}$ — the set of state variables it reads from ($D_{\text{in}}(e)$) and writes to ($D_{\text{out}}(e)$).

**Locality assumption (L).** For each event $e$, $|D(e)| < \infty$ and $D_{\text{out}}(e)$ is disjoint from $D_{\text{in}}(e')$ for all $e' \prec e$ not immediately preceding $e$. Events act on bounded local state domains.

**Causal dependency axiom (CD).** If event $e_2$ reads a state written by event $e_1$ — i.e., $D_{\text{out}}(e_1) \cap D_{\text{in}}(e_2) \neq \emptyset$ — then $e_1 \prec e_2$.

*Remark on CD.* This axiom links the abstract partial order relation $\prec$ (defined in Definition 1.2 structurally) to the physical notion of state propagation. Without it, $\prec$ is a purely formal ordering with no guaranteed connection to actual state domain dependencies. CD is the minimally necessary bridge: it says the causal ordering respects information flow. It is a physical assumption — not derivable from the poset axioms alone — and is stated explicitly here as such.
\end{definition}

\medskip

\section{Three Properties and Their
Derivation}\label{three-properties-and-their-derivation}

\subsection{P1: Causal Partial
Ordering}\label{p1-causal-partial-ordering}

\textbf{Claim.} \((E, \prec)\) is a strict partial order.

\begin{proof}
\end{proof}

\subsection{P2: Causal Invariance}\label{p2-causal-invariance}

\textbf{Claim.} Spacelike-separated events commute:
\(e_1 \perp e_2 \implies e_1 \circ e_2 = e_2 \circ e_1\).

\begin{proof}
\end{proof}

We claim \(D_{\text{out}}(e_1) \cap D_{\text{in}}(e_2) = \emptyset\) and
\(D_{\text{out}}(e_2) \cap D_{\text{in}}(e_1) = \emptyset\).

Suppose otherwise: say
\(s \in D_{\text{out}}(e_1) \cap D_{\text{in}}(e_2)\). Then by the
Causal Dependency Axiom (CD),
\(D_{\text{out}}(e_1) \cap D_{\text{in}}(e_2) \neq \emptyset\) implies
\(e_1 \prec e_2\) --- contradicting \(e_1 \perp e_2\). The symmetric
argument gives the second case.

Therefore \(D(e_1) \cap D(e_2) = \emptyset\): the events act on disjoint
state domains. Operators acting on disjoint domains commute (their
composition is independent of order). \qed

\textbf{Corollary.} P2 follows from P1, locality (L), and the Causal
Dependency Axiom (CD). The physical content is in CD --- the requirement
that the causal ordering respects state domain dependencies. CD is an
independent structural axiom, not derivable from the poset definition
alone.

\subsection{P3: Finite Minimum Event
Latency}\label{p3-finite-minimum-event-latency}

\begin{assumption}[There exists $\tau_{\min} > 0$ such that every causal event has latency $\tau_e \geq \tau_{\min}$.]
There exists $\tau_{\min} > 0$ such that every causal event has latency $\tau_e \geq \tau_{\min}$.

**Status: Postulate (well-motivated by local finiteness).** P3 is not derived from local finiteness alone — a genuine derivation would require showing that finite minimum latency follows from discrete poset structure without invoking the continuum limit, which would be circular (the continuum limit requires the discreteness of P3 to be defined). P3 is therefore treated as a well-motivated independent structural postulate, motivated by:
\end{assumption}

\begin{itemize}
\tightlist
\item
  Local finiteness: between any two causally related events, the causal
  chain is finite. This strongly motivates a minimum spacing but does
  not logically force a specific minimum time.
\item
  Physical calibration: the Planck time
  \(t_P = \sqrt{\hbar G / c^5} \approx 5.4 \times 10^{-44}\) s is the
  scale below which the classical spacetime description breaks down.
  Local finiteness is the structural claim; \(t_P\) is its physical
  calibration.
\end{itemize}

P3 joins local finiteness as the two structural postulates of the
cause-plex at Layer 0.

\medskip

\section{Time as Causal Path Count}\label{time-as-causal-path-count}

\begin{definition}[Causal path]
\label{definition:3-1}
A causal path from $e_1$ to $e_2$ is a sequence $e_1 = f_0 \prec f_1 \prec \cdots \prec f_n = e_2$ of events. The *length* of the path is $n$ (number of steps).

\end{definition}

\begin{definition}[Proper time]
\label{definition:3-2}
Given a reference cause-plex $\mathcal{C}_{\text{ref}}$ (an oscillating loop with period $T_{\text{ref}}$, e.g., an atomic clock), the proper time elapsed along a worldline $\gamma$ is:
\end{definition}

\[\tau(\gamma) = \frac{|\gamma|}{|\mathcal{C}_{\text{ref}}|}\]

where \(|\gamma|\) is the number of causal events along \(\gamma\) and
\(|\mathcal{C}_{\text{ref}}|\) is the number of events per reference
period. This is a dimensionless ratio until the reference period is
assigned a unit.

\textbf{This definition requires no background time coordinate.} Time is
a count of cause-plex events relative to a reference oscillation ---
intrinsic to the structure, not imposed from outside.

\textbf{Corollary (Q5 resolved).} Timescale separation between loops is
a ratio of causal event counts per reference period. Fast loops have
high event density per reference cycle; slow loops have low density.
Timescale is a derived quantity from cause-plex structure, not an
independent structural parameter.

\medskip

\section{Recovering the Lorentzian
Metric}\label{recovering-the-lorentzian-metric}

\subsection{Malament's Theorem}\label{malaments-theorem}

\textbf{Theorem (Malament 1977).} Let \((M, g)\) and \((M', g')\) be two
distinguishing spacetimes (spacetimes where distinct points have
distinct causal pasts and futures). If there exists a bijection
\(\phi: M \to M'\) that preserves the causal ordering in both directions
(a causal isomorphism), then \(\phi\) is a conformal isometry:
\(g' = \Omega^2 \phi^{\ast} g\) for some smooth positive function
\(\Omega\).

\textbf{Meaning.} The causal ordering uniquely determines the metric
\textbf{up to a conformal factor} \(\Omega^2\) --- an overall scale that
can vary point-by-point but cannot change the causal structure.
Malament's theorem means: know the causal partial order, know the metric
up to scale.

\subsection{Fixing the Conformal Factor by Event
Counting}\label{fixing-the-conformal-factor-by-event-counting}

\textbf{Claim.} The conformal factor \(\Omega\) is fixed by the density
of causal events per spacetime volume.

In causal set theory, the fundamental relation is:

\[N(A) \approx \frac{V(A)}{V_P}\]

where \(N(A)\) is the number of causal events in spacetime region \(A\),
\(V(A)\) is the spacetime volume of \(A\) in the continuum limit, and
\(V_P = \ell_P^4\) is the Planck volume (the natural unit set by local
finiteness).

This is the \textbf{``number = volume''} conjecture of causal set theory
(Bombelli et al.~1987, Sorkin 1991). If the event count faithfully
approximates spacetime volume, the conformal factor is fixed: regions
with higher event density have smaller \(V_P\)-normalized volume,
setting \(\Omega\) consistently across the manifold.

\textbf{Together:} causal ordering (Malament) + event counting (number =
volume) \(\to\) full Lorentzian metric. No additional assumptions beyond
P1 and local finiteness.

\subsection{\texorpdfstring{The Metric Signature
\((-,+,+,+)\)}{The Metric Signature (-,+,+,+)}}\label{the-metric-signature--}

The Lorentzian signature is not assumed --- it follows from the
asymmetry built into the causal partial order.

\textbf{Claim.} The causal partial order \((E, \prec)\) naturally
induces a bilinear form with signature \((-,+,+,+)\) in 3+1 dimensions.

\textbf{Argument.} In the continuum limit, consider the interval
function \$I(e\_1, e\_2) = \$ number of events causally between \(e_1\)
and \(e_2\):

\begin{itemize}
\tightlist
\item
  If \(e_1 \prec e_2\): \(I(e_1, e_2) > 0\) --- timelike separation,
  events share a causal cone
\item
  If \(e_1 \perp e_2\): \(I(e_1, e_2) = 0\) by definition --- spacelike
  separation, no causal path
\item
  On the boundary: \(I(e_1, e_2) = 0\) with \(e_1 \prec e_2\) ---
  lightlike, causal but no intermediate events
\end{itemize}

The interval function is the continuum spacetime interval \(ds^2\) up to
sign convention. The single negative direction (timelike) vs.~three
positive directions (spacelike) reflects the observed 3+1 dimensionality
of the cause-plex --- an empirical fact about the physical world's
cause-plex, not a derivation from pure structure.

\textbf{On dimensionality.} The specific signature \((-,+,+,+)\) vs
\((-,+,+)\) or \((-,+,+,+,+)\) is not derivable from the abstract
cause-plex structure alone --- it is a fact about how many independent
spacelike dimensions the physical cause-plex has. The structure
determines the form of the metric; the dimensionality is a property of
the specific cause-plex realized. See OP3.

\subsection{Lorentz Invariance}\label{lorentz-invariance}

\textbf{Claim.} In the continuum limit, the symmetry group preserving
causal structure is the Lorentz group.

\textbf{Argument (following Wolfram 2020, adapted).} A transformation
\(\Lambda\) of the cause-plex preserves causal structure if:

\[e_1 \prec e_2 \iff \Lambda(e_1) \prec \Lambda(e_2)\]

By Malament's theorem, causal-structure-preserving maps are conformal
isometries. The subgroup of conformal isometries that also preserve
event count density (fixing the conformal factor) is the isometry group
of the metric. For flat Minkowski spacetime, this is the Poincaré group.
The homogeneous subgroup (fixing the origin) is the Lorentz group.

In terms of causal path counts: a Lorentz boost is a transformation that
changes the distribution of causal events between timelike and spacelike
directions while preserving the total causal interval. Time dilation and
length contraction are the observed consequence of this redistribution.
\qed

\medskip

\section{Energy from Symmetry}\label{energy-from-symmetry}

\subsection{Discrete to Continuous
Symmetry}\label{discrete-to-continuous-symmetry}

The cause-plex has \textbf{discrete time-translation symmetry} if the
transition rules governing causal events are the same at every step ---
i.e., the cause-plex is homogeneous along causal chains. In the
continuum limit (\(\tau_{\min} \to 0\), many events per reference
period), discrete translation symmetry becomes continuous
time-translation symmetry.

\subsection{Noether's Theorem}\label{noethers-theorem}

\textbf{Theorem (Noether 1915).} For every continuous symmetry of the
action of a physical system, there is a corresponding conserved
quantity.

\textbf{Applied to time-translation symmetry:} continuous
time-translation invariance of the cause-plex action gives a conserved
quantity. We define this quantity as \textbf{energy}.

\textbf{Meaning.} Energy is not primitive --- it is the name for the
conserved quantity associated with the cause-plex having the same event
structure at every time step. In regions where this symmetry holds (most
of macroscopic physics), energy is well-defined and conserved. In
regions where it is broken (strongly non-equilibrium, rapidly evolving,
symmetry-breaking transitions), ``energy'' is not a clean quantity and
the framework works at the causal event level directly.

\textbf{Units.} Energy has units J = kg·m²·s⁻² where all three base
units are ratios of cause-plex path counts to reference oscillations: -
\textbf{Second:} \(9{,}192{,}631{,}770\) Cs-133 hyperfine transition
cycles (by definition) - \textbf{Meter:} distance light travels in
\(1/299{,}792{,}458\) seconds = \(c / f_{\text{ref}}\) cause-plex
propagation events - \textbf{Kilogram:} defined via Planck's constant
\(h = 6.626 \times 10^{-34}\) J·s, itself a count of action quanta
(phase cycles of matter waves)

All units reduce to cause-plex event counts relative to reference
oscillations. No unit is assumed.

Similarly, \textbf{spatial translation symmetry} \(\to\) momentum;
\textbf{rotational symmetry} \(\to\) angular momentum; \textbf{U(1)
gauge symmetry} \(\to\) charge.

\medskip

\section{The Coarse-Graining Ladder}\label{the-coarse-graining-ladder}

The cause-plex structure, once spacetime and conserved quantities are
established, produces the coarse-grained descriptions used throughout
the epimechanics series:

{\def\LTcaptype{none} % do not increment counter
\begin{longtable}[]{@{}
  >{\raggedright\arraybackslash}p{(\linewidth - 6\tabcolsep) * \real{0.2500}}
  >{\raggedright\arraybackslash}p{(\linewidth - 6\tabcolsep) * \real{0.2500}}
  >{\raggedright\arraybackslash}p{(\linewidth - 6\tabcolsep) * \real{0.2500}}
  >{\raggedright\arraybackslash}p{(\linewidth - 6\tabcolsep) * \real{0.2500}}@{}}
\toprule\noalign{}
\begin{minipage}[b]{\linewidth}\raggedright
Scale
\end{minipage} & \begin{minipage}[b]{\linewidth}\raggedright
Cause-plex description
\end{minipage} & \begin{minipage}[b]{\linewidth}\raggedright
Coarse-grained description
\end{minipage} & \begin{minipage}[b]{\linewidth}\raggedright
Valid where
\end{minipage} \\
\midrule\noalign{}
\endhead
\bottomrule\noalign{}
\endlastfoot
Planck & Discrete causal events \((E, \prec)\) & --- & Always \\
Quantum field & Gauge boson exchange events & Quantum fields
\(\hat{\phi}(x)\) & Flat spacetime, weak coupling \\
Classical mechanics & Many-event averages & Force \(F = dp/dt\), energy
\(W\) & \(\hbar \to 0\), macroscopic \\
Thermodynamics & Ensemble of causal event chains & Temperature, entropy
& Many-body, near-equilibrium \\
Biology & Coupled oscillating cause-plexes & ATP, metabolic flux,
\(\rho_{\text{ac}}\) & Time-translation symmetry holds locally \\
Institution & Human-scale cause-plexes & Dollars, decisions, norms &
Social-scale description valid \\
\end{longtable}
}

Each row is a coarse-graining of the row above. The descriptions in
column 3 are valid Layer C observables within their scope. ``Energy
exchange'' at the biological scale is valid because time-translation
symmetry holds there; it is not the primitive.

\medskip

\section{Curved Spacetime and General
Relativity}\label{curved-spacetime-and-general-relativity}

\subsection{\texorpdfstring{Non-Uniform Event Density \(\to\) Curved
Spacetime}{Non-Uniform Event Density \textbackslash to Curved Spacetime}}\label{non-uniform-event-density-to-curved-spacetime}

In flat Minkowski spacetime, causal event density is uniform across the
cause-plex. Curved spacetime corresponds to non-uniform event density:
regions with higher event density per coordinate volume have more causal
structure per unit coordinate interval --- the cause-plex is ``denser''
there.

The stress-energy tensor \(T_{\mu\nu}\) is a measure of energy-momentum
density at a point --- which is a measure of cause-plex event density
and flow at that point. The Einstein field equation:

\[G_{\mu\nu} = 8\pi T_{\mu\nu}\]

states that cause-plex event density (\(T_{\mu\nu}\)) determines the
curvature of the cause-plex metric (\(G_{\mu\nu}\)). This is a
cause-plex interpretation of GR.

\subsection{Actions in the Cause-Plex
Framework}\label{actions-in-the-cause-plex-framework}

Two distinct actions appear in the cause-plex framework, operating at
different layers. They should not be conflated.

\textbf{Gravitational action (Layer 0 --- purely combinatorial).}
Benincasa and Dowker (2010, \emph{Phys. Rev.~Lett.} 104:181301) define a
scalar curvature action on causal sets from event counts alone --- no
energy, no Lagrangian:

\[S_{\text{BD}} = \ell_P^2 \sum_{x \in \mathcal{C}} \left[ 1 - N_1(x) + \frac{N_2(x)}{2} - \cdots \right]\]

where \(N_k(x)\) counts causal intervals of depth \(k\) above \(x\). In
the continuum limit, \(S_{\text{BD}} \to \int R \sqrt{-g}\, d^4x\) ---
the Einstein-Hilbert action. This is clean at Layer 0: no energy
assumed, no circularity.

\textbf{Matter action (Layer C --- definitional, not derived).} The
matter action \(S[\gamma] = \sum_k \Delta E(e_k) \cdot \tau_{e_k}\) is a
Layer C construct. Energy \(\Delta E\) is defined at Layer C as the
Noether conserved quantity where time-translation symmetry holds
(Section 5). Time \(\tau\) is the event-count ratio (Definition 3.2). At
Layer C, the action is therefore well-defined and non-circular: it uses
quantities that are precisely defined at that layer.

The earlier concern (AUDIT.md Priority 1a) was that energy appeared to
be both derived-from and used-in the action. This concern is dissolved
by the layer separation: energy is not derived from the matter action at
Layer 0; rather, the matter action is a Layer C expression that names
what the Layer 0 phase function \(\phi(\gamma)\) (proved to lie in U(1)
in Complex Amplitudes from Loop-Phase
corresponds to in Layer C vocabulary. The identification
\(\phi(\gamma) = S[\gamma]/\hbar\) is a definitional bridge between
layers, not a circular derivation. See Proposition 6.2 of that paper for
the full argument.

\textbf{\(\hbar\) non-circularly.} The quantum of action
\(\hbar = \Delta E_{\min} \cdot \tau_{\min} = \Delta E_{\min}/|\mathcal{C}_{\text{ref}}|\)
--- minimum energy per causal event divided by the reference clock
frequency --- is defined using only Layer C energy and the event-count
time definition. It does not presuppose \(\hbar\). See §6.3 of the
loop-phase paper.

\subsection{Status: Partial (GR)}\label{status-partial-gr}

Deriving the Einstein field equation from cause-plex dynamics (rather
than interpreting it) requires showing that the dynamics governing
causal event density necessarily produce the GR equations in the
continuum limit. The Benincasa-Dowker action (Section 7.2) provides a
combinatorial causal-set action that converges to the Einstein-Hilbert
action --- this is the strongest existing result. Full derivation of the
dynamics remains open. This is the subject of:

\begin{itemize}
\tightlist
\item
  \textbf{Causal dynamical triangulations (CDT)} --- numerical evidence
  that the path integral over causal geometries gives 4D de Sitter
  spacetime in the continuum limit (Ambjorn, Jurkiewicz, Loll 2004)
\item
  \textbf{Causal set dynamics} (Sorkin's sequential growth models) ---
  show that certain natural random growth processes on causal sets
  produce Lorentzian geometry
\item
  \textbf{Spin foam models} --- discrete path integrals over causal
  structures that give GR in the semiclassical limit
\end{itemize}

\textbf{Current status:} The cause-plex interpretation of GR is
structurally correct and consistent with these results. A complete
derivation of the Einstein field equation from cause-plex primitives
alone is open. This is where epimechanics connects to the frontier of
quantum gravity research.

\medskip

\section{Relationship to Causal Set
Theory}\label{relationship-to-causal-set-theory}

At the physics level, the cause-plex framework is built on causal set
theory (CST). This section states the overlap and contributions
honestly.

\subsection{What the cause-plex shares with causal set
theory}\label{what-the-cause-plex-shares-with-causal-set-theory}

The cause-plex \((E, \prec)\) as defined in Section 1 is identical to a
causal set as introduced by Bombelli, Lee, Meyer, and Sorkin (1987): a
locally finite partially ordered set where the partial order encodes
causal precedence. The following results in this paper are drawn
directly from CST:

{\def\LTcaptype{none} % do not increment counter
\begin{longtable}[]{@{}
  >{\raggedright\arraybackslash}p{(\linewidth - 2\tabcolsep) * \real{0.5000}}
  >{\raggedright\arraybackslash}p{(\linewidth - 2\tabcolsep) * \real{0.5000}}@{}}
\toprule\noalign{}
\begin{minipage}[b]{\linewidth}\raggedright
Result
\end{minipage} & \begin{minipage}[b]{\linewidth}\raggedright
CST source
\end{minipage} \\
\midrule\noalign{}
\endhead
\bottomrule\noalign{}
\endlastfoot
P1: causal partial order & Bombelli et al.~(1987), Definition \\
P2: spacelike commutativity & Standard in CST; follows from locality \\
Metric from causal order (Malament) & Malament \cite{Malament1977},
applied in CST \\
Number = volume conjecture & Sorkin (1991), foundational CST
conjecture \\
Lorentz invariance from causal structure & Henson (2006), Dowker
(2006) \\
Gravitational action (Section 7.2) & Benincasa \& Dowker (2010) \\
QM via path integral on causal sets & Sorkin \cite{Sorkin1994} quantum
measure theory \\
\end{longtable}
}

The cause-plex framework does not re-derive these results --- it
inherits them. Any paper that presents these results as novel without
acknowledging CST would be misleading.

\subsection{What the cause-plex framework
adds}\label{what-the-cause-plex-framework-adds}

The distinctive contributions of epimechanics beyond CST are:

\textbf{1. The coarse-graining ladder to biology and society.} CST
focuses on the Planck-scale structure of spacetime. Epimechanics extends
the same primitive to bonds, loops, auto-causal entities, meta-entities,
and social systems --- a multi-scale program that CST does not address.

\textbf{2. The Q1--Q4 entity type descriptors.} The classification of
entities by bond type, loop order, basin depth, entropy production, and
repair rate {[}Part 1.5: Causors{]} is not present in CST. This provides
the vocabulary for the entity taxonomy.

\textbf{3. The selective convergence argument for dimensionality.} CST
takes the dimensionality of spacetime as an empirical input calibrated
by the number-volume conjecture. The cause-plex framework adds the
observer-selection argument (Section 10) grounding why 3+1 is the
observed dimensionality in terms of structural stability conditions for
observer-class entities.

\textbf{4. The loop-phase derivation of complex amplitudes.} The
argument that \(i\) in \(e^{iS/\hbar}\) follows from stable loop
structure {[}causeplex\_loop\_phase.md{]} is not present in CST.
Sorkin's quantum measure theory derives complex amplitudes from a
different direction (probability theory hierarchy).

\textbf{5. Integration with the epimechanics entity framework.} The
connection between fundamental causal structure and the
biological/social entity hierarchy is the central application of
epimechanics and has no analog in CST.

\subsection{Positioning}\label{positioning}

The cause-plex framework is best understood as: \emph{causal set theory
at the physics level, with a coarse-graining program extending to
biological and social systems, grounded by an observer-selection
argument for dimensionality, and augmented by a loop-phase derivation of
quantum amplitudes.} The physics is not new; the extension and
integration are.

\medskip

\section{Summary: What Is Derived, What Is
Assumed}\label{summary-what-is-derived-what-is-assumed}

{\def\LTcaptype{none} % do not increment counter
\begin{longtable}[]{@{}
  >{\raggedright\arraybackslash}p{(\linewidth - 4\tabcolsep) * \real{0.3333}}
  >{\raggedright\arraybackslash}p{(\linewidth - 4\tabcolsep) * \real{0.3333}}
  >{\raggedright\arraybackslash}p{(\linewidth - 4\tabcolsep) * \real{0.3333}}@{}}
\toprule\noalign{}
\begin{minipage}[b]{\linewidth}\raggedright
Claim
\end{minipage} & \begin{minipage}[b]{\linewidth}\raggedright
Status
\end{minipage} & \begin{minipage}[b]{\linewidth}\raggedright
Derivation
\end{minipage} \\
\midrule\noalign{}
\endhead
\bottomrule\noalign{}
\endlastfoot
Causal partial order exists (P1) & \(\checkmark\) Derived & By
construction from Definition 1.2 \\
Spacelike events commute (P2) & \(\checkmark\) Derived & From P1 +
locality (L) + Causal Dependency Axiom (CD) (Section 2) \\
Finite minimum event latency (P3) & \(\ast\) Postulate & Well-motivated
by local finiteness; not formally derived (Section 2) \\
Proper time as event count ratio & \(\checkmark\) Derived & Definition
3.2 \\
Q5 (timescale) is derived, not primitive & \(\checkmark\) Derived &
Corollary to Definition 3.2 \\
Metric determined up to conformal factor & \(\checkmark\) Derived &
Malament's theorem (1977) \\
Conformal factor fixed by event counting & \(\checkmark\) (conditional)
& Number = volume conjecture (well-supported, not proven) \\
Full Lorentzian metric falls out & \(\checkmark\) (conditional on above)
& Malament + event counting \\
Lorentz invariance & \(\checkmark\) Derived & Malament + continuum
limit \\
Metric signature \((-,+,+,+)\) & \(\triangle\) Observer-selection
argument & Causal structure gives \(n_t = 1\); \(n_s = 3\) follows from
stability filter (Tangherlini + knot topology) applied to observer-class
entities. This is a rigorous observer-selection argument, not a
derivation from pure cause-plex structure (Section 10) \\
Energy from time-translation symmetry & \(\checkmark\) Derived &
Noether's theorem in continuum limit \\
All units as event count ratios & \(\checkmark\) Derived & Reference
oscillation definitions \\
Flat spacetime (Minkowski) & \(\checkmark\) Derived & Uniform event
density \\
Curved spacetime interpretation & \(\checkmark\) Derived & Non-uniform
density = curvature \\
Einstein field equation & 🔬 Open & Frontier of causal set / quantum
gravity research \\
\end{longtable}
}

\textbf{One remaining empirical input:} the spatial dimensionality of
the physical cause-plex (3+1). The framework derives that the metric has
Lorentzian signature; the specific \((-, +, +, +)\) form requires that
the cause-plex has 3 independent spacelike dimensions. This is an
observed property of the physical world, not a logical necessity.

\medskip

\section{Why 3+1: Selective
Convergence}\label{why-31-selective-convergence}

The question of why we observe \((-,+,+,+)\) specifically --- not
\((-,+,+)\) or \((-,+,+,+,+)\) --- is addressed in the companion paper
Why 3+1: Observer-Selection Constraints on Spacetime

\textbf{Summary of the argument:} - \(n_t = 1\) follows directly from P1
(the causal partial order requires exactly one timelike direction;
\(n_t = 0\) breaks causal ordering, \(n_t \geq 2\) gives ultrahyperbolic
PDEs with no well-posed initial value problem). - \(n_s = 3\) follows
from two independent stability constraints applied to observer-class
entities: Tangherlini's orbital stability result (\(n_s > 3\) gives
unstable atoms) and the knot-topology requirement for topologically
non-trivial loop structures (\(n_s < 3\) makes all loops isotopic,
excluding \(\mathrm{CI} > \mathrm{CI}_{\min}\)).

This is an \textbf{observer-selection conjecture}, not a derivation from
the cause-plex primitive. It is honestly labeled as such throughout.

\medskip

\bibliographystyle{plain}
\bibliography{causeplex_spacetime}

\end{document}
